\section{Related Work}
\label{sec:related}

\para{Linear representations in LLMs.}
The linear representation hypothesis posits that high-level concepts are
encoded as directions in LLM activation spaces. \citet{marks2023geometry}
show that truth and falsehood are linearly separable in residual
streams, with simple difference-in-mean probes generalizing across
datasets. \citet{nanda2023emergent} find that Othello-playing networks
learn linear board-state representations that can be manipulated via
vector addition. \citet{park2023linear} formalize the hypothesis using
counterfactual semantics and introduce the causal inner product, proving
that causally separable concepts are orthogonal. \citet{park2024geometry}
extend this framework to categorical and hierarchical concepts, showing
that hierarchies are encoded as orthogonal subspaces. Our work builds
directly on these geometric predictions, testing whether the theoretical
orthogonality translates to practical compositionality.

\para{Function vector composition.}
\citet{todd2023function} discover that in-context-learned functions are
represented as compact vectors and test their composition via
parallelogram arithmetic. Composition succeeds for tasks with separable
sub-operations (\eg selection $+$ transformation) but fails for
entangled transformations (\eg translation). \citet{peng2025linear}
identify linear transformations between related knowledge in logit space
and show that composition requires both high correlation \emph{and} high
transformation precision. \citet{khandelwal2025compose} study two-hop
factual recall and find that models default to direct shortcuts when
linear mappings exist, bypassing compositional processing entirely.
Unlike these works, which study composition of \emph{functional}
directions, we study composition of \emph{concept} directions---testing
whether ``formal'' $+$ ``positive sentiment'' yields a representation
encoding both properties.

\para{Activation steering and its limitations.}
Contrastive activation addition~\citep{stolfo2024improving} applies
instruction-specific steering vectors to improve instruction following.
Critically, they find that summing multiple vectors at a single layer is
``largely unsuccessful,'' and instead apply vectors at different layers.
\citet{tan2024analyzing} rigorously evaluate steering vector reliability,
finding that per-sample steerability is bimodal---nearly 50\% of inputs
can be anti-steerable for some concepts. \citet{chalnev2024improving}
use sparse autoencoders to measure unintended side effects of steering
and propose SAE-targeted steering to minimize collateral activation.
These findings motivate our study: if individual steering vectors are
already unreliable, understanding when their \emph{composition} works is
essential for safe deployment.

\para{Compositionality in neural representations.}
Beyond steering, a broader literature examines how neural networks
represent compositional structure. \citet{guo2025quantifying} measure
additive compositionality using CCA and find that middle-to-late
transformer layers show the strongest compositional signals, declining at
the final layer. \citet{nagar2025transformer} test six composition models
across embedding architectures and find that ridge regression and vector
addition best account for compositional meaning. \citet{reif2025vocab}
show that morphological variations can be captured as additive
transformation vectors, succeeding for inflectional but failing for
derivational morphology. \citet{dziri2023faith} argue that transformers
solve compositional tasks via linearized subgraph matching rather than
systematic composition. Our work complements these studies by focusing
specifically on \emph{which} concept directions compose in the residual
stream and \emph{what geometric properties predict} composability.
